\scanpage{138}
\mspage{PU 50}

Le même terme peut être yang dans sa relation à tel autre, et yin dans sa relation à tel autre encore. Voir à ce sujet la section ``L'ambiguité créatrice (1) : paires, ribambelles et rondes'' (n° 3).

\noindent\rule{8em}{.4pt}

Comme je l'ai souligné dès le début de cette réflexion, il semble bien qu'une fois choisi l'ensemble $\mathrm T$ des ``termes'', représentant les qualités et entités cosmiques qu'on se propose d'étudier, la structure de graphe orienté correspondante est déterminé sans ambigu\leavevmode\rlap{i}\handwritten{\"{\phantom{\i}}}té. C'est-à-dire, pour deux termes $\mathrm a$ et $\mathrm b$ dans $\mathrm T$, on peut décider (par une intuition ou une réflexion de nature ``philosophique'', visiblement extra-mathématique) si ces deux termes ``forment couple'', et si oui, lequel des deux termes y joue le rôle yang (ou encore, doit figurer comme ``origine'' de l'arête orientée joignant les deux sommets représentatifs de $\mathrm a$ et de $\mathrm b$).

La deuxième structure qui est intervenue jusqu'à présent, se superposant à la première, est la structure d' a f f i n i t é. En langage mathématique courant, c'est encore une structure de g r a p h e (mais non orienté cette fois), consistant dans la prescription, parmi l'ensemble de toutes les ``paires'' possibles d'éléments de $\mathrm T$ (i.e. de parties de $\mathrm T$ réduites à deux éléments $\mathrm a,\ \mathrm b$), d'un certain sous-ensemble, formé des paires $\{\mathrm a,\mathrm b\}=\{\mathrm b,\mathrm a\}$ pour lesquelles $\mathrm a$ et $\mathrm b$ sont considéré\deleted{e}s comme étant ``voisins'', ou comme ``présentant des affinités''. Cette notion d'affinité est, elle aussi, de nature ``philosophique'', mais cette fois beaucoup moins clairement définie. Il n'y aura guère de lecteur qui ne décèlera, sans nuance d'une hésitation, un ``affinité'' entre ``rêve'' et ``imagination'', ou entre ``rêve'' et ``le possible''. Par contre, la question si ``rêve'' est voisin de ``hasard'', ou de ``jeu'', voire de ``liberté'' ou d' ``abandon'', aura sûrement des réponses bien différentes d'une personne à l'autre, et même chez la même personne, suivant les dispositions dans lesquelles elle abordera cette question. En fait, ce que nous révèle une intuition philosophique plus ou moins exercée, ce n'est pas tant une information du type ``tout ou rien''(``$\mathrm a$ et $\mathrm b$ sont voisins'', ou : ``ils n'ont rien à voir'') mais plutôt du type ``plus ou moins'' (comme : ``$\mathrm a$ et $\mathrm b$ sont très voisins'', ou ``assez voisins'', ou ``vaguement apparentés''...). C'est ce ``vague'' indissociable, semblet-t-il, de la notion d'affinité (dans le contexte qui nous intéresse), qui est cause aussi de l'arbitraire que je signalais dès le début, pour la formation des ``groupes'' de couples et des ``affinités'' entre tels groupes, par quoi former le diagramme (non orienté) des ``portes sur l'Univers'', alias ``diagramme en arbre de Noël''.

On est donc dans une situation un peu de guingois, où le mathématicien habitué à travailler avec des structures bien définies, se verrait confronté à une \margin{\handwritten{\unsure{$\times$|*}}}sorte de ``structure floue'', où il est censé tirer partie (\deleted{à} on ne\deleted{\unsure{l'|y}} sait trop encore
% archive sheet number: 138
